\begin{abstract}
Industrial visual inspection must recognize rare defects even when only a
handful of defect-free examples are available. This project studies few-shot
anomaly detection on six VisA categories. We first build a template-difference
baseline that compares each product image with its closest normal reference.
We then reproduce AnomalyDINO with a frozen DINOv2 ViT-S/14 backbone at
$448\times448$ resolution and agnostic preprocessing. A transparent
multi-layer improvement fuses normalized patch features from transformer blocks
9 and 11 before nearest-neighbor retrieval. All methods use the same official
1-class split, reference filenames, calibration images, and random seeds.
We evaluate image-level detection and patch-level localization over
1, 2, 4, and 8 shots with seeds 0, 1, and 2. The experiments show that frozen
foundation features provide a large advantage over raw template differences,
while combining a late semantic layer with an earlier spatial layer changes the
trade-off between anomaly detection and localization. The complete code,
configuration files, saved reference lists, metrics, visualizations, and
runnable demo are available at \url{\ProjectURL}.
\end{abstract}
